% ===============================
% MAT221 -- Real Analysis
% Riemann Integral
% ===============================

\documentclass[12pt,a4paper]{article}

% ===============================
% Metadata
% ===============================
\newcommand*{\CourseCode}{MAT221}
\newcommand*{\CourseTitle}{Real Analysis}
\newcommand*{\Chapter}{Lebesgue’s Criterion for Integrability}
\newcommand*{\Topics}{Measure Theory, Lebesgue’s Theorem}
\newcommand*{\DocDate}{September 06, 2026}

% ===============================
% Header
% ===============================
\input{header}

% ===============================
% Document
% ===============================
\begin{document}

\FrontPage
\Large


\begin{heading}
    Sets of Measure Zero
\end{heading}

\clearpage

\begin{theorem}
    A set $A \subseteq \mathbb{R}$ is said to be of measure zero if for every $\epsilon > 0$, there exists a countable collection of open intervals $\{U_n\}_{n=1}^{\infty}$ such that
    \[A \subseteq \bigcup_{n=1}^{\infty} U_n \text{ and }\sum_{n=1}^{\infty} \mu(U_n) < \epsilon,\] where $\mu(U_n)$ denotes the measure of the interval $U_n$.
\end{theorem}

\clearpage

\begin{theorem}[Countable Sets are of Measure Zero]
    Every countable set $A \subseteq \mathbb{R}$ is of measure zero.
\end{theorem}

\clearpage

\begin{heading}
    Continuity and Oscillation
\end{heading}

\begin{definition}[Oscillation of a Function]
    Let $f: [a, b] \to \mathbb{R}$ be a bounded function. The oscillation of $f$ at a point $x_0 \in [a, b]$ is defined as
    \[
    \omega_f(x_0)
    =
    \inf_{\delta>0}
    \left(
    \sup_{|x-x_0|<\delta} f(x)
    -
    \inf_{|x-x_0|<\delta} f(x)
    \right).
    \]
    
\end{definition}

\clearpage

\begin{theorem}
    Let $f: [a, b] \to \mathbb{R}$ be a bounded function. Then $f$ is continuous at a point $x_0 \in [a, b]$ if and only if the oscillation of $f$ at $x_0$ is zero, i.e., $\omega_f(x_0) = 0$, where \[
    \omega_f(x_0)
    =
    \inf_{\delta>0}
    \left(
    \sup_{|x-x_0|<\delta} f(x)
    -
    \inf_{|x-x_0|<\delta} f(x)
    \right)
    \]

    \end{theorem}


\clearpage

\begin{heading}
    Lebesgue’s Theorem
\end{heading}

\vspace{10pt}

\begin{theorem}
    A bounded function $f: [a, b] \to \mathbb{R}$ is Riemann integrable if and only if the set of points where $f$ is discontinuous has measure zero.
\end{theorem}

\EndPage

\end{document}